\icode{N21L-MCH-CA-L3-261026-002}
\nceeID[9]{2021山东}{16}
海鸥捕到外壳坚硬的鸟蛤（贝类动物）后，有时会飞到空中将它丢下，利用地面的冲击打碎硬壳。一只海鸥叼着质量\qty{0.1}{kg}的鸟蛤，在\(H =\) \,\qty{20}{m}的高度、以\(v_0 =\) \,\qty{15}{m/s}的水平速度飞行时，松开嘴巴让鸟蛤落到水平地面上。取重力加速度\(g =\) \,\qty{10}{m/s^2}，忽略空气阻力。
\begin{enumerate}
	\item 若鸟蛤与地面的碰撞时间\(\Delta t =\) \,\qty{0.005}{s}，弹起速度可忽略，求碰撞过程中鸟蛤受到的平均作用力的大小\(F\)；（碰撞过程中不计重力）
	\item 在海鸥飞行方向正下方的地面上，有一与地面平齐、长度\(L =\) \,\qty{6}{m}的岩石，以岩石左端为坐标原点，建立如图所示坐标系。若海鸥水平飞行的高度仍为\qty{20}{m}，速度大小在\(15\)\(\sim\)\qty{17}{m/s}之间，为保证鸟蛤一定能落到岩石上，求释放鸟蛤位置的\(x\)坐标范围。
\end{enumerate}
\begin{tikzpicture}[>=stealth,font=\scriptsize]
	\draw[->,blue] (-.6,0)--(4,0)node[below]{ \(x\)};
	\draw[->,blue] (0,-0.5)--(0,3)node[left]{ \(y\)};
	\draw[fill=gray!50] (0,0) rectangle(.8,-.1);
	\draw[<->,thin] (0,-.25)--node[below]{ \(L\)}(.8,-.25);
	\draw[dashed,thin] (.8,0)--(.8,-.5);
	\node at(3.5,2.6){ 海鸥};
	\node[blue] at(-.2,-.2){ \(O\)};
\end{tikzpicture}
\begin{answer}
	（1）\qty{500}{N}；（2）[\qty{34}{m},\qty{36}{m}]或(\qty{34}{m},\qty{36}{m})
\end{answer}